% Generated by scripts/fit_fig7.py; do not edit by hand.
\begin{table*}[t!]
\centering
\begin{tabular}{lccccc}\hline\hline
data & $A$ & $C$ & $D$ & $\chi^2/\mathrm{dof}$ & AICc \\ \hline
noiseless, $\epsilon=0.05$ & $-0.0165^{+0.0002}_{-0.0002}$ & $1.96^{+0.03}_{-0.03}$ & $0.5162^{+0.0002}_{-0.0002}$ & 0.25/6 & 11.05 \\
noisy, without knitting & $-0.022^{+0.004}_{-0.005}$ & $2.6^{+0.3}_{-0.4}$ & $0.521^{+0.005}_{-0.004}$ & 2.70/6 & 13.50 \\
noisy, with knitting & $-0.024^{+0.005}_{-0.01}$ & $1.9^{+0.4}_{-0.6}$ & $0.519^{+0.011}_{-0.004}$ & 2.04/6 & 12.84 \\ \hline\hline
\end{tabular}
\caption{Weighted truncated cosine fits to the same nine time points shown in Fig.~\ref{fig:knitting-time-evolution}. Here $A$, $C$, and $D$ are the amplitude, frequency, and offset in Eq.~\eqref{eq:trunc_fit}. Each entry gives the bootstrap median and central 68\% parametric-bootstrap interval.}
\label{tab:fig7-fits}
\end{table*}

Using the model-selection rule described in the text, the truncated cosine is the
primary model.  The fitted knitted-minus-un-knitted frequency difference is
$\Delta C=-0.692^{+0.534}_{-0.647}$
(68\% parametric-bootstrap interval).
